\documentclass[12pt,a4paper]{article}

% ------------------------------------------------
% Encoding and Language
% ------------------------------------------------
\usepackage[utf8]{inputenc}
\usepackage[T1]{fontenc}
\usepackage[french]{babel}

% ------------------------------------------------
% Layout & Typography
% ------------------------------------------------
\usepackage[a4paper,margin=2.5cm, headheight=15pt]{geometry}
\usepackage{setspace}
\usepackage{microtype} % Improves text justification and kerning
\onehalfspacing

% ------------------------------------------------
% Graphics, Tables, and Colors
% ------------------------------------------------
\usepackage[table]{xcolor} % Loaded first to avoid clashes with colortbl
\usepackage{graphicx}
\usepackage{float}
\usepackage{booktabs}
\usepackage{array}
\usepackage{longtable}
\usepackage{tabularx}

% ------------------------------------------------
% Math
% ------------------------------------------------
\usepackage{amsmath}
\usepackage{amssymb}

% ------------------------------------------------
% Professional Design & Formatting
% ------------------------------------------------
\usepackage{titlesec}
\usepackage{fancyhdr}
\usepackage{caption}
\usepackage[many]{tcolorbox} % Upgraded to 'many' for advanced styling
\usepackage{enumitem} % For cleaner lists


% ------------------------------------------------
% Colors
% ------------------------------------------------
\definecolor{primary}{RGB}{22,61,120}
\definecolor{secondary}{RGB}{52,89,149}
\definecolor{lightgray}{RGB}{245,245,245}
\definecolor{accent}{RGB}{240,244,250} % Soft blue for table rows and boxes

% ------------------------------------------------
% Hyperlinks (Loaded last)
% ------------------------------------------------
\usepackage{hyperref}
\hypersetup{
    colorlinks=true,
    linkcolor=primary,
    filecolor=secondary,
    urlcolor=primary,
    pdftitle={AgentBench-FR},
    pdfpagemode=FullScreen,
}

% ------------------------------------------------
% Headers / Footers
% ------------------------------------------------
\pagestyle{fancy}
\fancyhf{}

\fancyhead[L]{\textbf{\textcolor{primary}{AgentBench-FR}}}
\fancyhead[R]{\textbf{\textcolor{secondary}{AI Agent Framework Benchmark}}}
\fancyfoot[C]{\textbf{\textcolor{primary}{\thepage}}}

% Colored header rule
\renewcommand{\headrulewidth}{1pt}
\renewcommand{\headrule}{\hbox to\headwidth{\color{primary}\leaders\hrule height \headrulewidth\hfill}}

% ------------------------------------------------
% Section Style
% ------------------------------------------------
\titleformat{\section}
{\Large\bfseries\color{primary}}
{\thesection}{1em}{}[\vspace{-0.2cm}\color{primary}\titlerule]

\titleformat{\subsection}
{\large\bfseries\color{secondary}}
{\thesubsection}{1em}{}

% ------------------------------------------------
% Lists
% ------------------------------------------------
\setlist[itemize]{label=\textcolor{primary}{$\bullet$}, itemsep=0.3em, leftmargin=*}

% ------------------------------------------------
% Tables
% ------------------------------------------------
\renewcommand{\arraystretch}{1.4}
\arrayrulecolor{primary}
\captionsetup[table]{labelfont={bf,color=primary}, textfont={it,color=secondary}}
\captionsetup[figure]{labelfont={bf,color=primary}, textfont={it,color=secondary}}

% ------------------------------------------------
% Boxes
% ------------------------------------------------
\newtcolorbox{keypoint}{
  colback=accent,
  colframe=primary,
  boxrule=0pt,
  leftrule=4pt,
  arc=2mm,
  fontupper=\sffamily,
  top=3mm,
  bottom=3mm,
  left=4mm
}

% ------------------------------------------------
% Placeholder Figure
% ------------------------------------------------
\renewcommand{\placeholderfig}[2]{
\begin{figure}[H]
\centering
\begin{tcolorbox}[
  colback=white,
  colframe=secondary,
  boxrule=1pt,
  arc=3mm,
  width=0.85\textwidth,
  halign=center,
  valign=center,
  drop fuzzy shadow
]
\vspace{0.5cm}
\textbf{\Large\textcolor{primary}{#1}}
\vspace{0.8cm}\\
{\linespread{1.5}\selectfont #2 \par}
\vspace{0.5cm}
\end{tcolorbox}
\caption{#1}
\end{figure}
}

% ------------------------------------------------
% Title
% ------------------------------------------------

\title{}
\author{}
\date{}

\begin{document}

% =====================================================
% COVER PAGE
% =====================================================

\begin{titlepage}

\centering

\vspace*{1.5cm}

{\Huge\bfseries
\textcolor{primary}{AgentBench-FR}
}

\vspace{0.5cm}

{\Large\color{secondary}
Évaluation comparative des frameworks d'agents IA
}

\vspace{1.2cm}

\textcolor{primary}{\rule{0.8\linewidth}{1pt}}

\vspace{0.6cm}

{\Large\bfseries
LangChain \textbullet{} CrewAI \textbullet{} AutoGen \textbullet{} LlamaIndex
}

\vspace{0.6cm}

\textcolor{primary}{\rule{0.8\linewidth}{1pt}}

\vspace{2.5cm}

\begin{tcolorbox}[
    colback=white,
    colframe=primary,
    boxrule=1.5pt,
    arc=4mm,
    width=0.85\textwidth,
    center,
    top=0.8cm,
    bottom=0.8cm
]
\centering
\renewcommand{\arraystretch}{1.5}
\begin{tabular}{>{\bfseries\color{primary}}l l}
Projet : & Benchmark IA reproductible \\
Objectif : & Évaluation comparative contrôlée \\
Version : & Rapport consolidé \\
Créé par: & Adam El Madani - Saif Eddine Faten - Bilal Ma \\
Encadré par: & Pr. Tarik Hajji \\
Date : & \today
\end{tabular}
\end{tcolorbox}

\vfill

{\Large\bfseries\color{secondary}
Projet de recherche
}

\vspace{1.5cm}

\end{titlepage}

% =====================================================
% ABSTRACT
% =====================================================

\begin{abstract}
\vspace{1cm}
\noindent AgentBench-FR propose une méthodologie reproductible permettant
de comparer plusieurs frameworks d'agents IA dans un environnement
strictement contrôlé.

Le benchmark évalue LangChain, CrewAI, AutoGen et LlamaIndex
à travers une campagne de 120 tâches couvrant plusieurs catégories :
recherche d'information, raisonnement séquentiel,
récupération documentaire,
collaboration multi-agent
et robustesse face aux ambiguïtés.

L'étude harmonise le backend d'exécution,
les paramètres de génération,
la structure des entrées-sorties
et les ressources matérielles
afin d'isoler l'impact réel du framework.

Les métriques observées incluent
la complétion,
la latence,
la consommation de ressources,
l'utilisation de tokens
et la stabilité d'exécution.
\vspace{1cm}
\end{abstract}

\newpage
{\hypersetup{linkcolor=black} % Keep TOC links black for readability
\tableofcontents
}
\newpage

% =====================================================
% CONTEXT
% =====================================================

\section{Contexte et objectif}

\begin{keypoint}

\textbf{Objectif principal}

Construire un benchmark reproductible permettant
de comparer objectivement plusieurs frameworks
d'agents IA dans des conditions identiques.

\end{keypoint}

Les frameworks d'agents IA permettent
d'orchestrer des chaînes de raisonnement,
d'utiliser des outils externes
et de coordonner plusieurs agents.

Cependant,
les comparaisons directes restent difficiles,
car les différences observées peuvent provenir :

\begin{itemize}

\item du modèle utilisé

\item du backend d'exécution

\item du matériel

\item des paramètres de génération

\item du style d'orchestration

\end{itemize}

AgentBench-FR vise à réduire ces biais
par une méthodologie strictement contrôlée.

% =====================================================
% PROBLEMATIC
% =====================================================

\section{Sujet et problématique}

Le projet compare :

\begin{itemize}

\item LangChain

\item CrewAI

\item AutoGen

\item LlamaIndex

\end{itemize}

La problématique principale est :

\begin{keypoint}

Comment comparer objectivement
des frameworks d'agents IA
ayant des architectures différentes
tout en contrôlant
les variables externes ?

\end{keypoint}

Hypothèses :

\begin{itemize}

\item AutoGen devrait favoriser les scénarios multi-agents

\item CrewAI devrait simplifier les workflows par rôles

\item LangChain devrait offrir une meilleure modularité

\item LlamaIndex devrait exceller sur les tâches documentaires

\end{itemize}

% =====================================================
% FRAMEWORKS
% =====================================================

\section{Frameworks étudiés}

\subsection{LangChain}

Framework modulaire
orienté chaînes de prompts,
mémoire,
agents
et outils.

\subsection{CrewAI}

Architecture basée sur
la spécialisation des rôles.

\subsection{AutoGen}

Interaction conversationnelle
entre plusieurs agents.

\subsection{LlamaIndex}

Framework spécialisé
dans les pipelines documentaires
et les scénarios RAG.

% =====================================================
% ARCHITECTURE
% =====================================================

\section{Architecture du benchmark}

\placeholderfig
{Architecture AgentBench-FR}
{
Tâches

$\downarrow$

Adaptateurs Framework

$\downarrow$

Backend Ollama

$\downarrow$

Runner Benchmark

$\downarrow$

Collecte métriques

$\downarrow$

Analyse JSON / CSV
}

Le benchmark utilise
une couche d'adaptation commune
afin d'assurer
une standardisation complète.

% =====================================================
% TASKS
% =====================================================

\section{Jeux de tâches}

\rowcolors{2}{accent}{white}

\begin{table}[H]

\centering

\small

\begin{tabular}{lcccccc}

\toprule

\rowcolor{primary}
\textcolor{white}{\textbf{Catégorie}} &
\textcolor{white}{\textbf{N1}} &
\textcolor{white}{\textbf{N2}} &
\textcolor{white}{\textbf{N3}} &
\textcolor{white}{\textbf{N4}} &
\textcolor{white}{\textbf{N5}} &
\textcolor{white}{\textbf{Total}} \\

\midrule

Recherche information & 5 & 5 & 5 & 5 & 5 & 25 \\

Raisonnement séquentiel & 5 & 5 & 5 & 5 & 5 & 25 \\

Récupération documentaire & 5 & 5 & 5 & 5 & 5 & 25 \\

Collaboration multi-agent & 5 & 5 & 5 & 5 & 5 & 25 \\

Erreurs contrôlées & 5 & 5 & 5 & 5 & 5 & 25 \\

\bottomrule

\end{tabular}

\caption{Répartition des tâches}

\end{table}

Total :

\begin{center}

\Huge\bfseries\color{primary}
120 tâches

\end{center}

% =====================================================
% METHODOLOGY
% =====================================================

\section{Méthodologie}

\begin{keypoint}

Toutes les expériences utilisent :

Backend unique \quad+\quad Même modèle \quad+\quad Même température \quad+\quad Même timeout \quad+\quad Même machine

\end{keypoint}

\rowcolors{2}{accent}{white}

\begin{table}[H]

\centering

\begin{tabularx}{\linewidth}{lX}

\toprule

\rowcolor{primary}
\textcolor{white}{\textbf{Paramètre}} &
\textcolor{white}{\textbf{Valeur}} \\

\midrule

Backend & Ollama local \\

Température & 0 \\

Timeout & Fixe \\

Format sortie & Structuré JSON \\

Répétitions & Contrôlées \\

\bottomrule

\end{tabularx}

\caption{Paramètres expérimentaux}

\end{table}

% =====================================================
% METRICS
% =====================================================

\section{Métriques}

Mesures observées :

\begin{itemize}

\item Taux de complétion

\item Latence

\item CPU

\item RAM

\item Utilisation GPU

\item Consommation tokens

\item Robustesse

\item Nombre d'appels outils

\end{itemize}

Erreurs :

\begin{itemize}

\item TIMEOUT

\item TOOL\_ERROR

\item LOOP

\item PARSE\_ERROR

\item CONTEXT\_OVERFLOW

\item UNKNOWN

\end{itemize}

% =====================================================
% RESULTS
% =====================================================

\section{Résultats}

\rowcolors{2}{accent}{white}

\begin{table}[H]

\centering

\small

\begin{tabular}{lcccc}

\toprule

\rowcolor{primary}
\textcolor{white}{\textbf{Framework}} &
\textcolor{white}{\textbf{Compl.}} &
\textcolor{white}{\textbf{Latence}} &
\textcolor{white}{\textbf{CPU}} &
\textcolor{white}{\textbf{RAM}} \\

\midrule

LangChain & 1.000 & 4.848 & 0.0018 & 38.14 \\

CrewAI & 1.000 & 4.814 & 0.0024 & 38.25 \\

AutoGen & 1.000 & 4.829 & 0.0027 & 38.22 \\

LlamaIndex & 1.000 & 4.466 & 0.0037 & 38.31 \\

\bottomrule

\end{tabular}

\caption{Synthèse globale}

\end{table}

% =====================================================
% FINDINGS
% =====================================================

\section{Analyse}

\begin{keypoint}

Principales observations :

\begin{itemize}

\item AutoGen performant sur workflows collaboratifs

\item LlamaIndex plus efficace sur corpus documentaires

\item LangChain flexible et généraliste

\item CrewAI améliore la séparation logique des rôles

\end{itemize}

\end{keypoint}

Les différences observées
proviennent principalement
des stratégies d'orchestration.

% =====================================================
% VALIDITY
% =====================================================

\section{Menaces à la validité}

\subsection{Validité interne}

Les abstractions propres
à chaque framework
peuvent influencer
le comportement.

\subsection{Validité externe}

Les résultats locaux
peuvent différer
dans des environnements cloud.

\subsection{Validité de construction}

Les métriques quantitatives
ne capturent pas entièrement
la qualité du raisonnement.

% =====================================================
% CONCLUSION
% =====================================================

\section{Conclusion}

\begin{keypoint}

AgentBench-FR fournit
une méthodologie claire
et reproductible
pour comparer objectivement
les frameworks d'agents IA.

\end{keypoint}

Le benchmark harmonise :

\begin{itemize}

\item les tâches

\item le backend

\item les paramètres

\item les ressources

\item les métriques

\end{itemize}

Cette approche améliore
la comparabilité
et la reproductibilité.

% =====================================================
% APPENDIX
% =====================================================

\appendix

\section{Annexes}

\subsection{Organisation dépôt}

\begin{itemize}

\item tasks.yaml

\item prompts.yaml

\item adapters/

\item run\_all.py

\item analyze.py

\item cleanup\_logs.py

\end{itemize}

\subsection{Configuration locale}

Backend : \texttt{Ollama} \\
Mode : \texttt{Local} \\
Sortie : \texttt{JSON + CSV}

\end{document}